% 第6章 6.5.1：三点路径单位脉冲 f=e_1 的谱展开与重构。
% 系数 hat f=(1/sqrt3,1/sqrt2,1/sqrt6)；各项 hat f_k u_k：
% (1/3,1/3,1/3)，(1/2,0,-1/2)，(1/6,-1/3,1/6)；相加得 (1,0,0)。
\begin{tikzpicture}[font=\figurefont\footnotesize]
\pgfplotsset{every axis/.style={width=31mm,height=28mm,
  ybar,bar width=2.6mm,xtick={1,2,3},xticklabels={$1$,$2$,$3$},
  axis x line=bottom,axis y line=left,enlarge x limits=0.32,ymin=-0.5,ymax=1.05,
  ytick={-0.5,0,0.5,1},yticklabel style={/pgf/number format/fixed},
  x axis line style={draw=FigureStroke,line width=.6pt},
  y axis line style={draw=FigureStroke,line width=.6pt},
  tick style={draw=FigureStroke},tick label style={font=\figurefont\scriptsize,text=FigureInk},
  extra y ticks={0},extra y tick style={grid=major,major grid style={draw=FigureStroke,line width=.6pt},tick label style={opacity=0}},
  ymajorgrids,grid style={draw=FigureGrid,line width=.4pt},
  title style={font=\figurefont\footnotesize}}}
\begin{axis}[name=A,title={$\tfrac{1}{\sqrt3}\bm u_1$}]
\addplot[fill=FigureBlue,draw=FigureBlue] coordinates {(1,0.3333) (2,0.3333) (3,0.3333)};
\end{axis}
\node at ($(A.east)+(3.5mm,0)$) {$+$};
\begin{axis}[name=B,at={(A.east)},anchor=west,xshift=7mm,yticklabels={},title={$\tfrac{1}{\sqrt2}\bm u_2$}]
\addplot[fill=FigureGreen,draw=FigureGreen] coordinates {(1,0.5) (2,0) (3,-0.5)};
\end{axis}
\node at ($(B.east)+(3.5mm,0)$) {$+$};
\begin{axis}[name=C,at={(B.east)},anchor=west,xshift=7mm,yticklabels={},title={$\tfrac{1}{\sqrt6}\bm u_3$}]
\addplot[fill=FigureYellow,draw=FigureYellow] coordinates {(1,0.1667) (2,-0.3333) (3,0.1667)};
\end{axis}
\node at ($(C.east)+(3.5mm,0)$) {$=$};
\begin{axis}[at={(C.east)},anchor=west,xshift=7mm,yticklabels={},title={$\bm f=\bm e_1$}]
\addplot[fill=FigurePurple,draw=FigurePurple] coordinates {(1,1) (2,0) (3,0)};
\end{axis}
\end{tikzpicture}
